\noindent\textbf{Context and objective.}
Modern Linux hosts execute application workloads alongside container runtimes,
orchestration services and administrative processes. Although containers
separate workloads at user-space level, their processes still depend on the
same host kernel for execution, credentials, files, memory and isolation. This
makes the kernel a useful observation point for activity that crosses
application and container boundaries.

Application logs expose only the information selected by their developers and
may omit operations performed through libraries, helper programs or compromised
processes. Kernel observations can provide a more uniform view, but their
meaning depends on where an operation is observed. A system-call entry records
what a process requested, a security hook exposes an operation under mediation,
and a later transition or return value may establish what occurred. A useful
monitor must preserve these distinctions while retaining enough process and
resource context to relate the resulting evidence.

Collecting this information also introduces practical constraints. A host
produces continuous background activity, so transporting every available event
can increase CPU and memory consumption while obscuring the operations relevant
to an investigation. Moreover, an enterprise kernel may provide vendor
backports and data layouts that differ from the corresponding upstream release.
The monitoring design must therefore combine selective collection with bounded
processing and compatibility verified on the actual target system.

This thesis designs, implements and evaluates a node-local runtime security
monitor based on eBPF and a Go user-space runtime. The prototype selectively
collects host events, converts them into a common typed representation and
evaluates them through configurable policies and deterministic detectors. The
reference environment is Rocky Linux 8.10 with an enterprise Linux 4.18
kernel. BTF and CO-RE are used to adapt access to kernel data structures, while
explicit compatibility paths address target-specific differences. The
prototype observes and detects activity without modifying kernel decisions.

\medskip
\noindent\textbf{Architecture and contributions.}
The central contribution is an end-to-end monitoring pipeline that combines
host-security observation points with a stable binary event contract,
selective loading, early filtering and local detection. Logical events form the
public interface, while a probe registry maps them to the physical eBPF
programs and internal dependencies required for collection. This separation
allows several hooks or system-call variants to produce one event schema
without exposing implementation details to policies and detectors.

Kernel operations are observed at complementary stages. System-call entry
preserves the caller's request, Linux Security Module hooks expose security
mediation, scheduler and credential hooks report transitions that reached
kernel state, and return probes add the result visible at the observed
boundary. Direct producers construct an event in one invocation; paired
producers retain entry data per thread and event, then complete it at return.
Per-CPU workspaces keep large records off the eBPF stack, and indexed arguments
allow unavailable fields to be omitted without changing the meaning of later
values. Only the populated record prefix is submitted through the perf buffer.

The implemented hooks cover process lifecycle and executable transitions,
credentials and process control, file and memory activity, namespaces,
cgroups, modules and selected kernel-facing interfaces. Shared helpers provide event
initialisation, bounded serialisation and an optional exact effective-UID
filter that can reject out-of-scope observations before enrichment and
transport.

\enlargethispage{\baselineskip}
\begin{center}
    \centering
    \setlength{\fboxsep}{5pt}
    \begin{tabular}{c}
        \fbox{\parbox[c][0.72cm][c]{0.86\textwidth}{\centering
            \textbf{Kernel observation points}\\
            System calls, tracepoints, kernel functions and LSM hooks}} \\[0.2em]
        $\downarrow$ \\[-0.05em]
        \fbox{\parbox[c][0.72cm][c]{0.86\textwidth}{\centering
            \textbf{Shared eBPF collection}\\
            Direct and paired producers, UID admission and indexed serialisation}} \\[0.2em]
        $\downarrow$ \\[-0.05em]
        \fbox{\parbox[c][0.58cm][c]{0.86\textwidth}{\centering
            Populated binary record through per-CPU perf-event transport}} \\[0.2em]
        $\downarrow$ \\[-0.05em]
        \fbox{\parbox[c][0.72cm][c]{0.86\textwidth}{\centering
            \textbf{Go processing pipeline}\\
            Schema decoding, policy admission and point or collective detection}} \\[0.2em]
        $\downarrow$ \\[-0.05em]
        \fbox{\parbox[c][0.58cm][c]{0.86\textwidth}{\centering
            Normalised events, evidence-preserving alerts and separate diagnostics}}
    \end{tabular}
    \par\vspace{0.35em}
    Figure 1: End-to-end path from kernel observation to user-space results.
\end{center}


\noindent\textbf{User-space processing and detection.}
The Go runtime validates the configuration, resolves the selected logical
events and expands their probe dependencies before loading the eBPF object.
The decoder interprets the common context and indexed arguments through the
event registry, producing the typed representation shared by downstream
components. Policies define the admitted event domain and suppression rules;
detectors evaluate the resulting evidence and generate alerts.

Point detectors match conditions over one event. Collective detectors retain
bounded partial state for short ordered sequences and correlate observations
through stable process identity, immediate process relationships, local
device--inode identity, cgroup identity or configured combinations of these
dimensions. Alerts preserve the matching observations in order and may carry
MITRE ATT\&CK tactics and techniques as classification metadata. Normalised
events and alerts are written separately from operational diagnostics, so
runtime failures are not confused with monitored activity.

\medskip
\noindent\textbf{Experimental evaluation.}
The prototype was evaluated on a two-core Minikube node using a vulnerable 
CTF workload. The monitor, target and triggering client ran in
separate Pods on the same node, allowing the experiment to verify cross-Pod
visibility through their shared host kernel. Three benign runs completed
without an alert. In all three exploit runs, the client recovered the 
flag and the detector produced exactly one collective alert
containing four ordered observations: execution of a shell, execution of
\emph{cat}, resolution of the flag file and the successful \emph{open} result.

Performance was measured at the monitor-container cgroup boundary. With exact
UID admission enabled, mean CPU consumption was 0.46\% of one core, compared
with 5.73\% for the same monitoring configuration without the filter. The
observed reduction was approximately 92\%, and the filtered profile remained
below the project's 5\% acceptance threshold. Peak cgroup memory was
92.17~MiB with filtering and 110.21~MiB without it. Falco and Tetragon 
were evaluated in separate three-run campaigns on the same
Minikube testbed and against the same exploit, both with and without a UID
condition. In the filtered profiles, the prototype achieved the lowest mean CPU
consumption: 0.46\% of one core, compared with 1.88\% for Tetragon and 5.15\%
for Falco, demonstrating the effectiveness of its kernel-side UID filter.

\medskip
\noindent\textbf{Conclusion and boundaries.}
This work shows that useful runtime visibility depends not only on collecting
kernel events, but on preserving where each observation occurs and how it
relates to the surrounding activity. By distinguishing requests, mediation
points and outcomes, then correlating their evidence in user space, the
prototype turns heterogeneous kernel observations into interpretable security
alerts. This provides a concrete basis for extending node-level monitoring
without coupling detection logic to individual kernel hooks.
